
\chapter{El \taqsim}
\label{taqsim2}

El \taqsim (en plural, \taqasimb) puede considerarse como la máxima expresión de un \maqamb. Un músico de ud, además de conocer una serie de piezas del repertorio árabe, debe ser también capaz de improvisar un \taqsim en los principales \maqamaatb.

En este capítulo, veremos algunas de las particularidades del \taqsim como forma musical y analizaremos algunos ejemplos.

\section{Estructura y características}

Un \taqsim clásico responde al siguiente esquema:

\begin{itemize}
\item Una introducción breve que presenta el \maqam utilizado. Se emplean frases típicas de comienzo, que permiten al oyente reconocer el \maqam en que se va a desarrollar el \taqsimb.
\item Un conjunto de frases en las que se produce el desarrollo del \taqsimb. El \emph{sayr} del \maqam define la progresión a lo largo de estas frases.
\item Un cierre para el que también se utilizan clichés melódicos propios del \maqam de partida, que ha de ser también el de finalización.
\end{itemize}

Funcionalmente, el \taqsim se puede entender como una pieza dividida de la siguiente manera:

\begin{itemize}
\item Presentación del \maqamb.
\item Establecimiento del \maqamb. Se enfatiza su estructura y su tónica para que el \maqam quede instaurado como fundamento de la pieza.
\item Desarrollo melódico.
\item Despliegue del \emph{sayr}. Desarrollo de modulaciones según lo habitual para el \maqam en cuestión.
\item Regreso al \maqam original y cierre.
\end{itemize}

La duración de un \taqsim puede ser desde un minuto hasta más de diez, sin que existan limitaciones en cuanto a su rítmica o su extensión. El \taqsim es una exploración de un \maqam que podríamos decir que tiene como únicas limitaciones aquellas que imponga el propio \maqam en el que se encuadra, a través de su \emph{sayr} y de los otros elementos que lo constituyen.

Cada frase tiene una serie de características, y es el desarrollo de estas el que produce la <<narrativa>> del \taqsimb. Entre estas características, se pueden destacar las siguientes:

\begin{itemize}
\item \textbf{Nota predominante.} Cada frase tiene una o varias notas sobre las que gravita el resto de la melodía. Enfatizar ciertas notas permite, asimismo, modular entre \yins a partir de ellas.
\item \textbf{Dirección de movimiento melódico.} Puede ser ascendente o descendente. El movimiento de una frase ha de estar en relación con los de las frases que lo preceden o suceden. En general, el \taqsim arranca de la nota tónica del \maqamb, asciende y modula, y finalmente desciende al \maqam original, por lo que las primeras frases acostumbran a tener una dirección ascendente, mientras que las finales son descendentes.
\item \textbf{Amplitud.} Una frase puede limitarse a unas pocas notas, un \yins, o bien recorrer todo el \maqam.
\item \textbf{Ritmo.} El \taqsim como pieza global no tiene estructura rítmica, y esta puede variar a criterio del interprete en cada frase.
\item \textbf{Duración.} El \taqsim alterna frases cortas con otras más largas. En general, las primeras presentan algún motivo melódico, y este se extiende después en los fragmentos de más duración.
\item \textbf{Pausa entre frases.} El silencio es un elemento importante de un \taqsim, y deben incluirse espacios suficientes entre frases. El efecto dramático de los silencios es una característica fundamental del \taqsimb.
\item \textbf{Cierre.} Las frases se cierran con un remate, conocido como \emph{qafla}, que ha de crear sensación de reposo, y que suele situarse sobre la nota tónica o alguna de las notas relevantes del \emph{maqam}. De especial relevancia es el remate final del \taqsimb, que concluye siempre sobre la nota tónica del \maqam.
\end{itemize}

El improvisador de un \taqsim debe buscar que esa narrativa tenga una continuidad y un flujo adecuado, al tiempo que debe incluir el suficiente contraste entre las distintas frases como para aportar interés.

El empleo de los anteriores elementos varía según el estilo musical, según lo explicado en el apartado \ref{estilos}. Así, en el estilo árabe las pausas entre frases suelen ser más largas y los cierres más resolutivos (el que culmina el \taqsim suele tener un carácter en ocasiones casi épico), con el fin de lograr un efecto dramático. En el iraquí, sin embargo, las pausas son mucho más breves y las resoluciones más discretas, lo que contribuye al carácter más intimista que caracteriza a este estilo.

\section{Acerca de la improvisación}

Es un hecho que la improvisación no puede enseñarse de la misma manera que otros elementos de la interpretación musical. Es asunto lleno de elementos subjetivos, y en el que es el propio intérprete el que debe inferir parte de su mecanismo por sí mismo. Sin embargo, la reticencia o miedo que se tiene muchas veces a la hora de aprender a improvisar está motivada por una falsa mística, o por la idea equivocada de que la improvisación es puramente un talento innato, y no uno que puede desarrollarse y que posee unas reglas y  mecanismos bien definidos.

Todo el mundo es capaz de improvisar. De hecho, todo el mundo improvisa y crea cuando mantiene una conversación, porque en esta no hay un guión escrito y es en el momento mismo de hablar que se eligen las palabras y frases, las expresiones y gestos.

La música, siendo un lenguaje, sigue un proceso similar cuando de improvisación se trata. A partir de un conocimiento de este lenguaje y de un mensaje a transmitir, se crea un discurso en el momento mismo de ejecutarlo.

Pensemos en la siguiente situación, a la cual, con más o menos arte, cualquiera puede hacer frente: contar de memoria un cuento clásico a un niño. Partimos de una historia conocida por todos (sea, por ejemplo \emph{Caperucita roja}), de la cual se saben sus ideas principales (quienes son los protagonistas, el comienzo y final del cuento, las acciones principales que suceden, etc.), y sobre esto se completa el desarrollo de la historia. El conocimiento del idioma y la soltura que tenemos en este nos permite crear un conjunto de frases que <<arman>> el cuento en base a ese puñado de conceptos que son los que contienen su esencia.

En un \taqsim, las piezas conocidas las aportan el \maqam y las propias características de la música árabe (sus peculiaridades de desarrollo melódico, su sonoridad, etc). En ese marco, el interprete <<rellena>> el resto del contenido y con ello da forma a un \taqsim.

Es decir, que la improvisación no un proceso completamente libre en el que se crea una melodía a partir de la nada. En realidad, los elementos fundamentales ya están definidos, y no se trata más que de darles una forma tangible y una estructura, sirviéndose para ello de un conjunto de patrones conocidos (de la misma manera que empleamos palabras y frases hechas cuando narramos un cuento del que conocemos su argumento). Con la particularidad, eso sí, de que ese trabajo se ha de hacer de manera inmediata, según se toca el \taqsimb.

Se requiere, por tanto, conocer el \maqam y sus pormenores (notas que lo componen, modulaciones habituales, etc), saber cómo trasladar todo ello al ud, y tener la soltura suficiente como para poder llevarlo a cabo de manera espontánea. Lo primero se logra con el simple estudio, y lo último mediante la práctica y la escucha y análisis de otros \taqasimb.

Contar bien un cuento en el supuesto anterior depende, entre otros aspectos, de cómo sepamos construir las frases, del arte que se tenga a la hora de añadir información adicional para hacerlo más interesante, o de la expresividad con la que lo narremos. Es decir, del contenido que creemos sobre esas premisas iniciales y de cómo sea la interpretación.

En el \taqsimb, sucede algo similar. La calidad del \taqsim la define nuestro talento a la hora de <<componer>> las frases que lo forman (a partir de nuestro conocimiento del \maqam y del vocabulario del que disponemos, no desde cero), así como la elección acertada de la ornamentación, que es la que aporta la expresividad narrativa.

\subsection{El entrenamiento de la improvisación}


***

\subsection{La repetición}

La repetición es un aspecto fundamental de un \taqsim y se da a varios niveles diferentes.

En primer lugar, la repetición de notas es un componente imprescindible de un \taqsimb, ya que es la manera principal de enfatizar una de ellas, y esto es necesario para establecer la nota tónica del \maqam o desplazar el eje melódico a otra nota de relevancia, por ejemplo a la hora de modular a un nuevo \maqamb.

En segundo lugar, la repetición de frases es también imprescindible para asentar la narrativa del \taqsimb. Utilizando frases y desarrollos melódicos que ya han aparecido antes en otras partes del \taqsim, se logra dotar a este de una estructura, evitando así que parezca un mero conjunto de notas arbitrarias.

La repetición de frases puede ser:

\begin{compactitem}
\item Repetición literal de una frase. Puede darse inmediatamente (una frase breve repetida varias veces), o bien con una cierta separación. En este segundo caso, es habitual el empezar con el mismo fragmento una frase corta del \taqsim y una posterior más larga que desarrolle la idea con mayor extensión. La primera serviría a modo de resumen previo de la segunda.
\item Repetición de la misma secuencia de notas, con un distinto patrón rítmico.
\item Repetición del mismo patrón rítmico con una distinta secuencia de notas. Un caso particular de esto son las secuencias melódicos (véanse los apartados \ref{secuenciasmelodicas} y \ref{secuenciasmelodicasvocabulario})
\end{compactitem}

Como regla general, un \taqsim debe tener algunas <<ideas principales>> (rítmicas, melódicas, de ornamentación, etc), que se han de repetir y ser fácilmente identificables, y una colección de otras notas que sirvan de amalgama para conectar esas ideas que aparecen a lo largo del \taqsimb, y que serán las que se encarguen de darle su identidad y significado.


\section{Análisis de algunos \taqsimb}

La mejor manera de comprender la idea de \taqsim para poder más tarde improvisar uno propio es escuchando y analizando un buen número de ellos. Esta sección incluye una descripción de algunos \taqasim en los \maqamaat más habituales, detallando aspectos como su estructura, la forma de sus frases o el \emph{sayr} seguido en cada caso.

Comencemos con los \taqasim que ya conocemos de un capítulo anterior, y de los que entonces analizamos su \emph{sayr}.

\subsection{\emph{Taqsim} Rast}

En primer lugar, un \taqsim en \maqam Rast.

\audioclip{\emph{Taqsim} Rast}

Podemos dividir este \taqsim en las siguientes partes:

\begin{itemize}
\item \textbf{Sección 1 (0:00 - 0:25)}: Presentación del \maqam. Frase introductoria breve. Énfasis sobre la nota tónica y la octava inferior, para dejar desde el inicio el \maqam que va a desarrollarse.
\item \textbf{Sección 2 (0:00 - 1:00)}: Primera frase. Nahawand en sol con dirección ascendente.
\item \textbf{Sección 3 (0:00 - 1:00)}: Final de la primera frase, que en este caso cubre los dos \ashnas del \maqam Rast (Rast en do y Rast en sol), y anticipa la siguiente sección.
\item \textbf{Sección 4 (0:00 - 1:00)}: Frase breve que arranca en Rast en sol y desciende a Rast en do. Presenta la frase que sigue, tanto en su estructura como en su modo melódico. Esta es una técnica habitual en el \taqsimb: presentar un motivo con una frase y después repetirlo de una manera más prolija y variada. La repetición a nivel estructural es un elemento fundamental de un \taqsimb.
\item \textbf{Sección 5 (0:00 - 1:00)}: Comienzo de la última frase que toma la estructura de la anterior y la desarrolla. Énfasis en do sobre la primera cuerda y el \yins Rast en sol.
\item \textbf{Sección 6 (0:00 - 1:00)}: Modulación a \yins Hiyaz en sol (es decir, \maqam Suznak). Antes de entrar en él, hay un par de notas que pertenecen a Nahawand en sol. Regreso a Rast en sol en la última parte.
\item \textbf{Cierre (0:00 - 1:00)}: Recorre todo el \maqam Rast en dirección descendente, con uso intenso del trémolo para añadir intensidad. Nótese como la frase final (\emph{qafla}) ya anticipa desde antes de su término la conclusión del \taqsim.
\end{itemize}

****
\subsection{\emph{Taqsim} Bayati}

\subsection{\emph{Taqsim} Nahawand}

\subsection{\emph{Taqsim} Rast}

A continuación, un nuevo \taqsim en \maqam Rast. Tanto este como los que siguen no se han visto en secciones anteriores, por lo que se añade también el análisis del \emph{sayr} y las modulaciones.

\subsection{\emph{Taqsim} Hiyaz}


\subsection{\emph{Taqsim} Ajam}